\section{Method}

\subsection{Overview}

Our Linformer implementation has four phases:

\begin{enumerate}
\item \textbf{Projection} ( 4.2): standard $Q, K, V$ projections.
\item \textbf{RoPE} ( 4.2): optional rotary embeddings.
\item \textbf{GQA expansion} ( 4.3): expand KV heads.
\item \textbf{Adaptive Linformer projection} ( 4.4): project $K, V$ to dimension $k$.
\item \textbf{Attention} ( 4.5): standard scaled-dot-product on projected tensors.
\item \textbf{Causal handling} ( 4.6): fallback to dense if needed.

\end{enumerate}
\subsection{Projections and RoPE}

Standard projections via LoRALinear:

\begin{lstlisting}
q = self.q\_proj(hidden\_states)
k = self.k\_proj(hidden\_states)
v = self.v\_proj(hidden\_states)
\end{lstlisting}

Reshape to multi-head:

\begin{lstlisting}
q = q.view(B, N, H, d\_h).transpose(1, 2)
k = k.view(B, N, H\_kv, d\_h).transpose(1, 2)
v = v.view(B, N, H\_kv, d\_h).transpose(1, 2)
\end{lstlisting}

Apply RoPE if position embeddings available.

\subsection{GQA Expansion}

\begin{lstlisting}
if self.kv\_repeat \textgreater{} 1:
    k = repeat\_kv(k, self.kv\_repeat)
    v = repeat\_kv(v, self.kv\_repeat)
\end{lstlisting}

After this, $K, V \in \mathbb{R}^{B \times H \times L \times d_h}$.

\subsection{Adaptive Linformer Projection}

\subsubsection{Three Cases}

Given sequence length $L$, projection dim $k$, and pre-allocated $L_\text{max}$:

\begin{table}[H]
\centering
\small
\begin{tabular}{lll}
\toprule
\textbf{Case} & \textbf{Condition} & \textbf{Action} \\
\midrule
\textbf{Short} & $L \leq k$ & Use $E[:L, :L]$, no compression \\
\textbf{Medium} & $L \leq L_\text{max}$ & Use $E[:k, :L]$, standard Linformer \\
\textbf{Long} & $L > L_\text{max}$ & Interpolate $E$ to length $L$ \\
\bottomrule
\end{tabular}
\end{table}

\subsubsection{Short Case ($L \leq k$)}

For short sequences, no compression is applied:

\begin{lstlisting}
E\_act = self.E[:L, :L]
K\_proj = E\_act @ k  # (B, H, L, d\_h)
\end{lstlisting}

This preserves exact attention for short sequences, ensuring no accuracy 
loss when compression isn't needed.

\subsubsection{Medium Case ($L \leq L_\text{max}$)}

Standard Linformer:

\begin{lstlisting}
k\_eff = min(self.k\_proj\_dim, L)
E\_act = self.E[:k\_eff, :L]
K\_proj = E\_act @ k  # (B, H, k\_eff, d\_h)
\end{lstlisting}

$K$ and $V$ are projected from $L$ to $k_\text{eff}$ positions.

\subsubsection{Long Case ($L > L_\text{max}$)}

For sequences longer than the pre-allocated max:

\begin{lstlisting}
E\_act = F.interpolate(
    self.E[:k\_eff].unsqueeze(0),
    size=L,
    mode="linear",
    align\_corners=False,
).squeeze(0)
\end{lstlisting}

Linear interpolation extends $E$ to length $L$. This is a graceful 
degradation -- not part of the original Linformer paper, but necessary for 
robustness.

\subsection{Attention on Projected Tensors}

\subsubsection{Scores}

\[
\text{scores} = \frac{Q \tilde{K}^\top}{\sqrt{d\_h}} \in \mathbb{R}^{B \times H \times L \times k\_\text{eff}}
\tag{4.1}
\]

Computed in FP32 for stability.

\subsubsection{Softmax}

Standard softmax over $k_\text{eff}$ dimension.

\subsubsection{Context}

\[
\text{context} = \text{attn} \cdot \tilde{V} \in \mathbb{R}^{B \times H \times L \times d\_h}
\tag{4.2}
\]

\subsubsection{Output Projection}

\begin{lstlisting}
context = context.transpose(1, 2).contiguous().view(B, L, self.all\_head\_size)
out = self.out\_proj(context * self.output\_scale)
\end{lstlisting}

Where \texttt{output\_scale} is a learnable scalar initialized to 1.

\subsection{Causal Handling for Decoders}

\textbf{The problem.} For a decoder, projecting $K$ via $E$ mixes future and past 
information:

\[
\tilde{K}\_j = \sum\_{i=1}^L E\_{j, i} K\_i
\tag{4.3}
\]

The resulting $\tilde{K}_j$ depends on \textbf{all} positions, not just $j$. A 
causal mask over $\tilde{K}_j$ is meaningless.

\textbf{The solution.} Fall back to exact dense attention:

\begin{lstlisting}
if self.is\_decoder and L \textgreater{} self.k\_proj\_dim:
    # Cannot compress without breaking causality
    causal\_scores = q @ k.transpose(-2, -1) / sqrt(d\_h)
    causal\_mask = tril(ones(L, L))
    scores = causal\_scores.masked\_fill(~causal\_mask, -inf)
    V\_proj = v  # Use unprojected V
\end{lstlisting}

This preserves causality at the cost of $\mathcal{O}(L^2)$ for decoders 
with $L > k$.

For short sequences ($L \leq k$), the projection is identity-like, so 
causality is approximately preserved. We use the projected path in this 
case.

\subsection{Full Algorithm}

\textbf{Algorithm 1:} Linformer forward pass
\begin{lstlisting}
Input:  X $\in$ R^{B$\times$L$\times$d}, k, L\_max
Output: Y $\in$ R^{B$\times$L$\times$d}

 1: Q $\leftarrow$ X W\_Q; K $\leftarrow$ X W\_K; V $\leftarrow$ X W\_V
 2: if RoPE: Q, K $\leftarrow$ RoPE(Q, K)
 3: if GQA: K, V $\leftarrow$ repeat\_kv(K, V)
 4: if L $\le$ k:
 5:     E\_act $\leftarrow$ E[:L, :L]; F\_act $\leftarrow$ F[:L, :L]
 6:     V\_proj $\leftarrow$ F\_act @ V
 7: elif L $\le$ L\_max:
 8:     k\_eff $\leftarrow$ min(k, L)
 9:     E\_act $\leftarrow$ E[:k\_eff, :L]; F\_act $\leftarrow$ F[:k\_eff, :L]
10:     V\_proj $\leftarrow$ F\_act @ V
11: else:
12:     E\_act $\leftarrow$ interpolate(E, size=L); F\_act $\leftarrow$ interpolate(F, size=L)
13:     V\_proj $\leftarrow$ F\_act @ V
14: K\_proj $\leftarrow$ E\_act @ K
15: scores $\leftarrow$ Q @ K\_proj^T / $\sqrt{}$d\_h
16: if is\_decoder and L \textgreater{} k:
17:     # Fallback: dense causal attention
18:     scores $\leftarrow$ Q @ K^T / $\sqrt{}$d\_h
19:     V\_proj $\leftarrow$ V
20:     scores $\leftarrow$ mask\_causal(scores)
21: scores $\leftarrow$ mask\_padding(scores)
22: attn $\leftarrow$ softmax(scores)
23: out $\leftarrow$ attn @ V\_proj
24: out $\leftarrow$ out\_proj(out * output\_scale)
25: return out
\end{lstlisting}

\subsection{Complexity Analysis}

\subsubsection{Encoder Complexity}

\begin{table}[H]
\centering
\small
\begin{tabular}{ll}
\toprule
\textbf{Operation} & \textbf{Complexity} \\
\midrule
$Q, K, V$ projections & $\mathcal{O}(L d^2)$ \\
$E K$ & $\mathcal{O}(k L d_h)$ \\
$Q \tilde{K}^\top$ & $\mathcal{O}(L k d_h)$ \\
Softmax & $\mathcal{O}(L k)$ \\
$\text{attn} \cdot \tilde{V}$ & $\mathcal{O}(L k d_h)$ \\
Output projection & $\mathcal{O}(L d^2)$ \\
\textbf{Total} & $\mathcal{O}(L k d_h + L d^2)$ \\
\bottomrule
\end{tabular}
\end{table}

For fixed $k$ and $d$: \textbf{linear in $L$}.

\subsubsection{Decoder Complexity}

\begin{table}[H]
\centering
\small
\begin{tabular}{ll}
\toprule
\textbf{Case} & \textbf{Complexity} \\
\midrule
$L \leq k$ & $\mathcal{O}(L k d_h)$ \\
$L > k$ & $\mathcal{O}(L^2 d_h)$ (fallback) \\
\bottomrule
\end{tabular}
\end{table}

\textbf{Trade-off.} Encoders benefit from linearity; long-sequence decoders 
fall back to quadratic.

\subsubsection{Memory}

\begin{table}[H]
\centering
\small
\begin{tabular}{ll}
\toprule
\textbf{Tensor} & \textbf{Size} \\
\midrule
$E, F$ & $2 k L_\text{max}$ \\
Attention scores & $\mathcal{O}(L k)$ \\
Context & $\mathcal{O}(L d_h)$ \\
\bottomrule
\end{tabular}
\end{table}

Much smaller than baseline $\mathcal{O}(L^2)$ for the scores matrix.

\subsection{Parameter Analysis}

\begin{table}[H]
\centering
\small
\begin{tabular}{llll}
\toprule
\textbf{Component} & \textbf{Baseline} & \textbf{Ours} & \textbf{Change} \\
\midrule
$W_Q, W_K, W_V, W_O$ & $4 d^2$ (frozen) & $4 d^2$ (frozen) & 0 \\
LoRA adapters & 0 & $4 r_L (d + d)$ & +small \\
$E, F$ & 0 & $2 k L_\text{max}$ & +medium \\
\texttt{output\_scale} & 0 & 1 & ~0 \\
\bottomrule
\end{tabular}
\end{table}

For $d = 768$, $k = 64$, $L_\text{max} = 512$:

\begin{table}[H]
\centering
\small
\begin{tabular}{ll}
\toprule
\textbf{Component} & \textbf{Params} \\
\midrule
Baseline per layer & 2,359,296 \\
LoRA (rank 16) & 98,304 \\
E, F & $2 \cdot 64 \cdot 512 = 65,536$ \\
output\_scale & 1 \\
\textbf{Total added} & ~163,841 \\
\textbf{Change} & \textbf{+6.9\%} \\
\bottomrule
\end{tabular}
\end{table}

We state this \textbf{explicitly} in the abstract.

\subsection{Theoretical Notes}

\subsubsection{Low-Rank Approximation Bound}

If the attention matrix $P$ has effective rank $r \ll L$, then:

\[
\|P - \tilde{P}\|\_F \leq \sqrt{\sum\_{i=r+1}^L \sigma\_i^2}
\tag{4.4}
\]

where $\sigma_i$ are singular values of $P$. For $k \geq r$:

\[
\|P - \tilde{P}\|\_F \approx 0
\tag{4.5}
\]

\subsubsection{Causal Preservation}

\textbf{Claim.} For decoders with $L > k$, the fallback preserves causality.

\textbf{Proof.} The fallback uses standard dense attention with a triangular mask, 
which is causally correct by construction. $\square$

\subsubsection{Identity Initialization}

\textbf{Claim.} With $E = I$, the projection $\tilde{K} = K$ and Linformer 
reduces to standard attention.

\textbf{Proof.} $E K = I K = K$. $\square$

This means our initialization starts at exactly the standard attention and 
deviates only as $E$ is trained.
\begin{lstlisting}

---
\end{lstlisting}